\chapter*{Kurzfassung}
\addcontentsline{toc}{chapter}{Kurzfassung}
\markboth{Kurzfassung}{Kurzfassung}
\thispagestyle{plain}

An gestapelten mobilen Bauzaunfeldern sind die Stirnflächen der Rahmenrohre zu
detektieren: kleine, nahezu identische Kreisringe, dicht gepackt und
wechselseitig verdeckt. Annotierte Datensätze existieren nicht, die
Handannotation hunderter Kleinobjekte pro Bild ist teuer. Der Ansatz umgeht es:
Trainiert wird \emph{ausschließlich} auf synthetischen Bildern aus neun
freigestellten Objektausschnitten und sieben Hintergrundaufnahmen, validiert
und getestet auf echten Aufnahmen.

Den Schwerpunkt trägt damit die \textbf{Augmentierungs- und Synthese-Pipeline},
nicht die Netzarchitektur: Sie erzeugt die Trainingsverteilung vollständig. Aus
den Grundlagen der Datenaugmentierung -- Kodierung von Invarianzen,
Rauschregularisierung, Mittel gegen \enquote{Shortcut Learning} -- entwickelt
die Arbeit eine Begründungspflicht: Jeder Parameter stützt sich auf eine an den
Realdaten \emph{gemessene} Eigenschaft, eine \emph{physikalische
Randbedingung} oder ein \emph{Literaturargument} -- sonst gilt er ausdrücklich
als \emph{Annahme}. Die Pipeline umfasst 36 Transformationen mit 43
Parameterwerten (6000 Bilder, 136\,384 gelabelte Objekte): 9 gemessen, 8
literaturgestützt, 3 physikalisch begründet, 9 gemischt, \textbf{14 als Annahme
gekennzeichnet}. Auf zwei realen Aufnahmen erreicht das Verfahren \mAPfifty{}
von \textbf{\zrefValMap} und \textbf{\zrefTestMap} bei einem besten $F_1$ von
\zrefValFf{} und \zrefTestFf; \code{test} berührt Training und Modellauswahl
nie. Der Sim-to-Real-Abstand zum synthetischen Split beträgt nur
\zsimRealVal{} bzw.\ \zsimRealTest\PP.

Geprüft wird die Pipeline an ihrer aufwendigsten Stufe, der Verdeckung:
\textbf{woraus die Verdecker bestehen dürfen}. Die aus Realaufnahmen
geschnittene Okkluder-Bibliothek wird durch zwei prozedurale ersetzt, eine
kalibrierte und eine naive -- gleicher Seed, byte-identische Labels, also ein
gepaarter Versuch mit allein variierenden Verdeckerpixeln. Ein naiv
gezeichneter grauer Balken kostet auf dem unberührten Testbild
\textbf{\zablBTestGap\PP{} \mAPfifty} und \zablBTestGapFf\PP{} $F_1$ -- fast
vollständig beim kleinsten Objektdrittel, wo der Rückruf von
\zBZeroTestRecKlein{} auf \zBTwoTestRecKlein{} fällt. Das an fünf Kennwerten
kalibrierte Zylindermodell erreicht dagegen dieselbe Leistung wie die
photographische Bibliothek: Entscheidend ist nicht die \emph{Herkunft} der
Verdeckerpixel, sondern ob ihre Statistik zur Umgebung passt.

Der größte Einzeleffekt ist schließlich keine Augmentierungsentscheidung,
sondern die Inferenzauflösung: \mAPfifty{} auf \code{val} steigt von
\zswValNiedrig{} bei 640\px{} auf \zswValBest{} (\zswValGain\PP). Und ein
systematischer Größenversatz von \zboxValOff{} bis \zboxTestOff\px{} zwischen
Generatorlabel und Handannotation erklärt den niedrigen \mAPseven{}-Wert auf
dem Testbild: Eine Sensitivitätsrechnung hebt ihn von \zsensTestMapSOrig{} auf
\zsensTestMapSNeu{}, während \mAPfifty{} und $F_1$ unberührt bleiben. Bei
programmatisch erzeugten Labels wird die Labelkonvention selbst zur messbaren
Größe.
